\documentclass[a4paper,12pt]{article}
\usepackage[utf8]{inputenc}
\usepackage[T1]{fontenc}
\usepackage[italian]{babel}
\usepackage{listings}
\usepackage{xcolor}
\usepackage{hyperref}
\usepackage[margin=3cm]{geometry}

\title{Documentazione Progetto Sistemi Operativi: PandOSsh Fase 3}
\author{Luca Argentino, Milo Disalvatore, Matteo Mazzetti, Zeno Maccari}
\begin{document}


\maketitle
\section{Gestione Memoria Virtuale (vmSupport.c)}
Il modulo \texttt{vmSupport.c} implementa il supporto per la memoria virtuale, la gestione dei TLB-Refill e il paging.

\vspace{0.5cm}
\subsection{Strutture Dati Globali}
Le variabili globali e le strutture definite nel modulo sono:
\begin{itemize}
    \item \texttt{swapPool}: Array di tipo \texttt{swap\_t} e di dimensione \texttt{POOLSIZE} (pari a \texttt{UPROCMAX * 2}, ovvero 16 frame) che rappresenta la tabella della Swap Pool per allocare le pagine in memoria fisica.
    \item \texttt{swapSemaphore}: Semaforo utilizzato per garantire la mutua esclusione della Swap Pool tra i vari processi.
    \item \texttt{holdingSwapMutex}: Cella di un'array che tiene traccia, tramite indice basato sull'ASID, di quale processo possiede attualmente il mutex per lo swap.
\end{itemize}

\vspace{0.5cm}
\subsection{Descrizione delle Funzioni}

\subsubsection{\texttt{void initPagetable(pteEntry\_t *pageTable, int asid)}}


\subsubsection{\texttt{void uTLB\_RefillHandler(void)}}


\subsubsection{\texttt{void initSwapPool()}}


\subsubsection{\texttt{int getFlashBlock(int vpn)} e \texttt{devreg\_t *getFlashRegister(int asid)}}
\begin{enumerate}
    \item \texttt{getFlashBlock}: Funzione ausiliaria che calcola il blocco corrispondente a un dato VPN. Restituisce 31 per l'indirizzo di stack (0xBFFFF), mentre per gli altri VPN sottrae l'offset base 0x80000.
    \item \texttt{getFlashRegister}: Calcola l'indirizzo fisico dei registri del Flash Device di un processo in base al suo ASID. Se il device non è valido ritorna \texttt{NULL}, altrimenti restituisce il puntatore alla struttura dei registri calcolata tramite l'offset della linea degli interrupt.
\end{enumerate}

\subsubsection{\texttt{int readWriteFlashdrive(int asid, int vpn, int phisicalFrame, int op)}}
\begin{enumerate}
    \item Ricava il numero di blocco logico dal VPN e il registro corretto del Flash Device associato al processo tramite il suo ASID.
    \item Scrive l'indirizzo del frame fisico iniziale  all'interno del campo \texttt{DATA0} del device.
    \item Costruisce il comando unendo il numero del blocco (shiftato di 8 bit) con il codice operativo (\texttt{op}), per poi inviare una richiesta sincrona sul campo \texttt{COMMAND} tramite una \texttt{SYSCALL(DOIO)}.
\end{enumerate}

\subsubsection{\texttt{void pager()}}
\begin{enumerate}
    \item Recupera il puntatore alla struttura di supporto del processo tramite \texttt{SYSCALL(GETSUPPORTPTR)} e verifica la causa del TLB fault dall'eccezione salvata.
    \item Se l'eccezione è di tipo \texttt{TLB-Mod} (modifica), il comportamento viene dirottato al gestore delle Program Trap.
    \item Acquisisce il mutex sulla Swap Pool tramite \texttt{PASSEREN} su \texttt{swapSemaphore} e aggiorna il tracker locale \texttt{holdingSwapMutex}.
    \item Sceglie un frame dalla Swap Pool utilizzando un algoritmo FIFO di tipo round-robin per rimpiazzare una pagina. Prima però controlla che ciò sia necessario: se è presente un frame non occupato utilizziamo quello.
    \item Se il frame è occupato da una pagina valida, la invalida rimuovendo il flag \texttt{VALIDON} dalla sua Page Table, svuota il TLB e salva il vecchio contenuto nel rispettivo Flash Device tramite \texttt{FLASHWRITE}.
    \item Legge la nuova pagina dal Flash Device del processo corrente e la carica nel frame fisico scelto tramite \texttt{FLASHREAD}. In caso di errore durante l'I/O, rilascia il semaforo e genera una Program Trap.
    \item Aggiorna l'entry della Swap Pool salvando il nuovo ASID, il numero di pagina (\texttt{missingPageNo}) e ricavando l'indirizzo della rispettiva entry nella Page Table.
    \item Modifica l'entry della Page Table appena trovata unendo l'indirizzo fisico del frame e impostando i flag \texttt{DIRTYON} e \texttt{VALIDON}.
    \item Rilascia l'uso del semaforo su \texttt{swapSemaphore} (resettando anche \texttt{holdingSwapMutex}) e restituisce il controllo ricaricando lo stato per tentare nuovamente l'istruzione \texttt{LDST}.
\end{enumerate}
\end{document}
